\begin{table}[htbp]
\centering
\caption{Respuestas Agregadas a un Shock Negativo de 10\% a la PTF en Manufactura}
\label{tab:mfg_agg}
\begin{threeparttable}
\begin{tabular}{@{}l c c c@{}}
\toprule
Variable & Impacto & Mín/Máx & Trimestre \\
\midrule
PIB & -4.862 & -4.862 & 1 \\
Inflación IPC & +4.427 & +4.427 & 1 \\
Consumo & -1.301 & -1.301 & 1 \\
Empleo & +2.286 & +2.286 & 1 \\
Tipo de Cambio Real & -2.121 & -2.121 & 1 \\
Balanza Comercial & -3.591 & -3.591 & 1 \\
Tasa de Política & +2.838 & +2.838 & 1 \\
\bottomrule
\end{tabular}
\begin{tablenotes}[flushleft]
\footnotesize
\item \textit{Notas:} PIB, consumo, empleo y TCR en \% de desviación del estado estacionario; inflación y tasa de política en pp anualizados; balanza comercial en pp del PIB. El shock es una reducción única de 10\% en la productividad total de factores del sector Manufactura con persistencia $\rho = 0.446$. Impacto = respuesta en el trimestre 1. Mín/Máx = respuesta extrema en 40 trimestres.
\end{tablenotes}
\end{threeparttable}
\end{table}
